% ========================================================================================
% APPENDICE - COMPONENTI ELEMENTARI RIUSATI
% ========================================================================================
\chapter{Componenti riusati}

In questa appendice sono raccolti i componenti elementari usati in più esercizi.

% ----------------------------------------------------------------------------------------
\section{Multiplexer 4:1}\label{app:mux4}
% ----------------------------------------------------------------------------------------

\progetto
Multiplexer a quattro ingressi descritto in stile dataflow con un'assegnazione
\texttt{with\ldots select}. È usato come blocco base del mux 16:1 (Esercizio~1.1).

\begin{porte}[Interfaccia dell'entity \ent{mux4}]
	\porta{d}{in}{4}{ingressi dati}
	\porta{s}{in}{2}{selezione}
	\porta{y}{out}{1}{uscita, $y = d_s$}
\end{porte}

\implementazione
\vhdlfile{comuni/mux4.vhd}{Multiplexer 4:1 (\texttt{mux4.vhd}).}{lst:mux4}

% ----------------------------------------------------------------------------------------
\section{Debouncer}\label{app:debouncer}
% ----------------------------------------------------------------------------------------

\progetto
Sincronizzatore a due flip-flop seguito da un filtro antirimbalzo: il nuovo livello del
pulsante viene accettato solo se resta stabile per \SI{10}{\milli\second}. Genera anche un
impulso lungo un ciclo di clock in corrispondenza della pressione.

\implementazione
\todo{\texttt{\textbackslash vhdlfile\{comuni/debouncer.vhd\}\{\ldots\}\{lst:debouncer\}}}

% ----------------------------------------------------------------------------------------
\section{Altri componenti}
% ----------------------------------------------------------------------------------------
\todo{Full adder, registro, contatore modulo N, ROM, RAM, driver del display a 7 segmenti.}
